\documentclass{article}

% CS230 project milestone template, built on the official NeurIPS 2026 style file.
% It uses the same setup as proposal_template/ and final_report_template/ so the
% proposal, milestone, and final report share one look. Keep the [final] option: it
% prints the author block and omits line numbers.
% natbib is loaded by the style file; this makes citations numeric, e.g. [1, 2].
% Remove the next line for author-year citations, e.g. (Vaswani et al., 2017).
\PassOptionsToPackage{numbers,compress}{natbib}
\usepackage[final]{neurips_2026}

% to avoid loading the natbib package, add option nonatbib:
%    \usepackage[final,nonatbib]{neurips_2026}

\usepackage[utf8]{inputenc} % allow utf-8 input
\usepackage[T1]{fontenc}    % use 8-bit T1 fonts
\usepackage{hyperref}       % hyperlinks
\usepackage{url}            % simple URL typesetting
\usepackage{booktabs}       % professional-quality tables
\usepackage{amsfonts}       % blackboard math symbols
\usepackage{nicefrac}       % compact symbols for 1/2, etc.
\usepackage{microtype}      % microtypography
\usepackage{xcolor}         % colors
\usepackage{graphicx}       % figures
% Page limit for the milestone, used in the text below. Update it here only.
\newcommand{\pagelimit}{3}

% Submission: upload this PDF on Gradescope.

% State the full title of your project. Keep the same title across the milestone and
% the final report where possible.
\title{Formatting instructions for CS230 project milestones}

% List the full name and SUNet ID of every team member.
% Use \And between authors to let LaTeX decide where to break the line, and
% \AND to force a line break at that point.
\author{%
  First Author \\
  \texttt{sunetid1@stanford.edu} \\
  \And
  Second Author \\
  \texttt{sunetid2@stanford.edu} \\
  %% examples of more authors
  %% \And
  %% Third Author \\
  %% \texttt{email} \\
  %% \AND
  %% Fourth Author \\
  %% \texttt{email} \\
}

\begin{document}

\maketitle

\section{Introduction}
The milestone is limited to \pagelimit~pages, excluding references. Write it as an
early draft of your final report so that you can reuse most of the text later. The
intended audience is the course staff, so do not spend space explaining standard
material. Explain the problem and why it is important or interesting, and be explicit
about the input and output: ``The input to our algorithm is an \{image, text prompt,
audio clip, patient record, video, etc.\}. We then use a \{CNN, transformer, diffusion
model, graph neural network, etc.\} to output a predicted \{label, caption, segmentation
mask, next action, etc.\}.'' Your final report must be on the same topic as this
milestone, so if the project has changed since your proposal, say what changed and why.

\section{Dataset}
Describe your dataset: where it comes from, roughly how many training, validation, and
test examples you have, and its license if it has one. Cite the dataset if it is from
prior work, such as ImageNet \cite{deng2009imagenet}. Say what you
have actually downloaded, collected, cleaned, or labeled so far, and what preprocessing
and data augmentation you have implemented, such as resizing or normalizing images,
tokenizing text, resampling audio, or discretizing time series. If you build on
pretrained encoders or embeddings, say which ones and what they are applied to. You may
show a few samples from your data here.

\section{Approach}
Describe the steps you have completed so far. If you are implementing an algorithm, you
should have started the implementation. We would like to see a model description and a training strategy (loss function for instance). Include
the relevant notation; formulas from the lectures are fine. Name the baseline you are
comparing against, and cite the methods you build on with \verb+\cite+, e.g.\ ResNet
\cite{he2016deep} or the Transformer \cite{vaswani2017attention}.

\paragraph{Remaining work.}
Describe precisely the remaining work you expect to complete before the final report,
and in what order, so that the TAs can give you feedback on the plan. Be brief: a short
paragraph or a few bullet points is enough.

\section{Partial Results}
Report whatever you have run so far, even if preliminary: the data you have collected and
processed, the performance of the baseline you have implemented and evaluated, and
preliminary results of your own method, if any. State your primary metric and report it
on held-out data rather than the training set. A small table such as
Table~\ref{tab:results} or a single figure is enough at this stage; see the commented-out
block below for how to include a figure. If a result is surprising or an assumption did
not hold, say so briefly and explain what you will change.

\begin{table}[h]
  \centering
  \begin{tabular}{lc}
    \toprule
    Method & Validation accuracy \\
    \midrule
    Baseline   & 0.71 \\
    Our method & 0.74 \\
    \bottomrule
  \end{tabular}
  \caption{Example results table. Replace with your own numbers and metric.}
  \label{tab:results}
\end{table}

% \begin{figure}[h]
%   \centering
%   \includegraphics[width=0.6\linewidth]{training_curve.pdf}
%   \caption{Training and validation loss per epoch for the baseline.}
%   \label{fig:curve}
% \end{figure}

\section*{References}
References do not count toward the page limit.
Include the papers you cite above, the dataset, and any code or libraries you build on.
% Acceptable formats
% include MLA, APA, and IEEE; any style is fine as long as it is consistent and each entry
% has author(s), title, venue, publisher, and year.

% The list below is generated by BibTeX from references.bib; run
% "latexmk -pdf main.tex" (or pdflatex, bibtex, pdflatex, pdflatex).
\renewcommand{\bibsection}{}   % the "References" heading is already typeset above
\medskip
{\small
\bibliographystyle{plainnat}
\bibliography{references}
}

\end{document}
